Lần sau cùng -- vâng, hóa ra đó là lần sau cùng -- tôi gặp anh là ở một buổi giới thiệu sách giáo khoa của nhóm Cánh Buồm. Có rất nhiều các anh chị quen biết đến dự hôm ấy và ban đầu, do đến muộn, tôi đã không kịp nhìn thấy anh. Khi phần trình bày của nhóm Cánh Buồm chấm dứt, tôi định ra về thì mới gặp anh. Đã biết về các vấn đề sức khỏe của anh trong thời gian sau này, còn nhớ rõ một lần đến thăm anh ở nhà và nghe chị giải thích về bệnh tình của anh, hôm ấy gặp và thấy anh vui, khỏe đi tham dự sinh hoạt tôi rất mừng. Anh bảo: Mình đi ăn trưa với nhau đi. Hai anh em đón taxi đi đến quán Ngon, anh chọn. Tình cờ được dịp gặp anh, lại được đi ăn với anh là một may mắn vượt quá mong ước của tôi. Lắm lần về Hà Nội vì làm việc túi bụi tôi không dám cả hẹn đến thăm anh chị. Rất may với cái duyên hôm ấy.

Biết tình trạng sức khỏe của anh, phải thú thật, ngồi với anh bữa ăn trưa hôm ấy, lòng dạ tôi luôn có một nỗi xót xa, đến độ ăn gì, món ăn ngon dở ra sao tôi chẳng còn có thể quan tâm. Trong lúc ấy, anh vẫn vui vẻ. Anh chọn món ăn, anh giới thiệu với tôi về các món anh chọn. Vẫn là anh, thân tình, dung dị như ông anh mà tôi quen biết đã hơn ba mươi năm. Nhưng, đàng sau tất cả những gì quen thuộc ấy, đã rõ là anh yếu đi nhiều, sự tinh nhanh của anh cũng suy giãm đi nhiều. Và tôi không thể nào không bận tâm, lo lắng về sự thay đổi ấy. Giờ nhớ lại, hôm ấy có lúc tôi như hồn vía ở đâu đâu.

Còn nhớ, trong lúc chờ thức ăn, anh đã làm một việc rất... anh. Anh bảo, anh gọi điện thoại cho chị để cho biết anh đang ở đâu. Anh cho chị biết anh đang đi ăn với tôi. Và tôi nhờ anh cho tôi nhắn lòi chào chị. Tôi nói, “rất anh” vì trong thâm tâm, từ khi biết anh chị, cách đối xử của hai anh chị với nhau vẫn là một ấn tượng đẹp đối với tôi. Mà anh chị thì xem bọn tôi như những đứa em nên những gì chúng tôi được chứng kiến là rất thật, rất tự nhiên. Anh Diệu, người chồng, người cha là một khía cạnh về anh mà tôi luôn trân trọng.

*
Các anh bên tạp chí Diễn Đàn nhắn, nếu tôi có viết được gì về anh thì gửi sang để các anh ấy đăng. Tôi đắn đo mãi. Chuyện về anh, công việc chuyên ngành của anh ở Việt Nam thì các anh chị từng có dịp cộng sự với anh sẽ viết tốt hơn, đúng và rõ hơn tôi nhiều. Còn về những suy nghĩ của anh về đất nước thì đã có những bài viết của anh về hiện tình hay hướng đi tương lai cần chọn cho đất nước, có lẽ tốt hơn hết là đăng lại những gì anh đã nghĩ suy và gởi gấm (lại). Những điều anh viết, phát biểu, dẫu đã nhiều năm trước, quý báu thay -- mà cũng xót xa thay -- giờ vẫn còn chưa cũ, giờ vẫn còn là những thôi thúc. Cho nên, trong những ngày này, trong nỗi buồn đưa tiễn anh, tôi chỉ có thể viết xuống những ghi nhớ rời và riêng tư về anh. Những ghi nhớ giờ đã thành hoài niệm của tôi. Về anh, một ông anh lớn, một người mà tôi đã thầm cố gắng học đôi điều, để rồi giờ đây phải nhận ra, học theo anh là một việc không dễ dàng, dù chỉ đôi điều tôi dám xem là vừa sức.

Anh là một nhà khoa học hệ thống, có nếp tư duy hệ thống. Và chính anh đã dùng trí tuệ đó để soi sáng cô đọng một vấn đề liên quan đến nghề nghiêp của tôi, cho tôi, ngay từ năm 1984. Dạo đó, tôi hành nghề trong lĩnh vực hệ thống máy tính nhằm phục vụ cho ứng dụng quản lí. 
Nên khi về thăm Viện Khoa học Tính toán và Điều khiển của anh, tôi đã không tránh khỏi điều quan tâm thật sơ đẳng: làm sao để Việt Nam quan tâm và đầu tư nhiều hơn đến việc áp dụng tin học (thời thuật ngữ CNTT còn chưa xuất hiện) vào quản lí kinh tế, xã hội. Anh trả lời ngắn gọn, đại ý: để làm điều đó, người ta phải quản lí bằng những con số (dữ liệu) thật. Khi người ta chưa quản lí bằng những con số thật thì máy tính chưa thể đắc dụng.

Anh nói đơn giản, bằng thứ ngôn từ so ra rất đơn sơ về một vướng mắc, một mâu thuẫn rất căn cơ, để giúp tôi, một người ở “ngoài” có thể nắm được cái khó của việc áp dụng các tri thức và công nghệ tiên tiến vào hoàn cảnh Việt Nam. Tùy trình độ của tôi, tùy khả năng chiêm nghiệm của tôi, câu trả lời ấy dẫn tôi đi qua khá nhiều giai đọan khác nhau của tiến bộ ngành, đối chiếu với tác động của ngành vào thực tế Việt Nam. Bây giờ là 34 năm sau, tôi vừa đọc trên Facebook câu chuyện quản lí bằng “hai sổ”. Điều anh nói, đơn giản về hình thức, vẫn không cũ, bất kể ta nhìn sự việc ở tầm bao la hay chi li.

Nhìn thấy như vậy, hiểu sâu sắc hiện trạng như vậy, nhưng anh vẫn luôn nhiệt tình, tích cực đóng góp cho sự nghiệp vận dụng công nghệ thông tin và máy tính vào phát triển xã hội. Đó cũng là một bài học quí báu cho tôi. Những năm có dịp trao đổi với anh về Chương trình Công nghệ Thông tin Quốc gia với anh, tôi đã được nghe một số phân tích, dẫn giải của anh. Tôi đã có dịp nhìn thấy một số khó khăn anh phải đương đầu, nhưng sự quan tâm và nhiệt tình của anh thì không vì thế mà suy suyễn.

Người ta vẫn hay nhắc đến những va chạm anh đã gặp phải trong cuộc sống và sự nghiệp chuyên ngành ở Việt Nam. Tuy ở xa, qua các nguồn tin, tôi vẫn được nghe phong thanh về những chuyện ấy, Nhưng, gặp anh, dù ở Hà Nội hay ở California, anh vẫn không hề bỏ phí thời gian để tỉ tê về các tin đồn ấy. Đôi lúc, phải chạm đến nó, anh lại dùng một cách nói vui để nói về những điều không vui ấy. Chỉ phớt qua thôi.

*
Vào những năm cuối 1980, đầu 1990, dư luận cũng có những bàn tán về một số bài viết, phát biểu của anh liên quan đến hiện tình đất nước và các biện pháp quản lí đất nước của giới cầm quyền. Thậm chí, có lúc tình hình dẫn đến những lo ngại về an ninh bản thân của riêng anh. Về nước, đến thăm anh chị, tất nhiên chúng tôi cũng lo lắng hỏi thăm anh. Lúc ấy, chị là người trả lời. Chị nói đại ý: Sợ chứ chú. Nên mỗi lần anh Diệu đi đâu chị đều phải đi theo. Đề phòng có gì xảy ra cho anh. Cách chị nói làm tôi xúc động lắm. Vì tôi nhìn chị, người đàn bà “tay không tấc sắt” ấy chỉ có vũ khí tình yêu và lòng quả cảm cá nhân để vung lên khi cần, hầu may ra có thể che chở cho anh, trong một tình huống nào đó có nhiều khả năng là không đơn giản, không thể xem nhẹ. Những lần có dịp về Hà Nội và anh chị bảo đến “tá túc” ở nhà anh chị, tôi càng hiểu về vai trò của chị trong cuộc sống của anh. Và với tình cảnh hôm nay, càng thấy cảm thương chị biết chừng nào.

Thế nhưng, ngay ở những lúc khó khăn như vậy, anh vẫn có cách hành xử của anh. Anh không hề phản ứng như một người “chống đối”, anh không chọn thái độ “đối lập” có vẻ phổ biến trong giai đoạn ấy. Khi người ta nghi ngại, rình rập anh, anh đã sẵn sàng “mở” ra cho họ cơ hội kiểm soát, ghé mũi vào những gì họ muốn săm soi, kiểm soát. Chính anh đã minh bạch với họ, trấn an họ trước. Anh đã cho tôi biết như thế, vì anh muốn tôi hiểu rằng, các cách trao đổi thư tín đều sẽ được chính anh “mở” ra cho những thế lực cần rình rập anh. Đàng khác, về một phía khác, trong một câu chuyện riêng, anh kể tôi nghe, giọng kể có chút vui đùa: đã có người trách anh sao không mạnh dạn lên tiếng chống đối chính sách đàn áp những người không cùng chính kiến với chế độ? Và anh kể luôn cách anh trả lời sự trách cứ ấy. Anh hỏi lại người bạn kia: Nếu các anh không bị cấm đoán, các anh có quyền nói thì các anh sẽ nói điều gì? Nói gì cho có lợi cho đất nước? Cách đặt vấn đề ấy, theo tôi, giải thích (cho tôi) cách anh chọn để truyền đạt các ý kiến của mình.

Anh đã kiên nhẫn, thẳng thắn và hoà nhã chọn dịp để trình bày những suy nghĩ của anh cho những người có thẩm quyền, ở vị trí (có thể) làm quyết định, để (may ra) họ chịu lắng nghe, chịu hiểu và (biết đâu) họ sẽ làm theo. Và anh đã làm đúng như vậy. Làm đúng phần vụ và trách nhiệm của một trí thức tâm huyết với đất nước như vậy. Nói vui, thái độ của anh, dưới mắt tôi, là một thái độ “sĩ phu” đích thực, dù anh không bao giờ chọn cho mình một thứ nhãn hiệu nào.

Và chủ trương ấy của anh rất nhất quán. Lần sau cùng đón anh ở nam Cali, tôi và một số bè bạn đã có buổi họp mặt rất vui với anh. Trong số có những người làm báo chí, có quan hệ nhiều với truyền thông. Thế là có đề nghị phỏng vấn anh về tình hình đất nước và các suy nghĩ liên quan của anh. Nhưng anh đã khéo léo và dứt khoát từ chối. Anh không muốn “người ta” gán cho anh cái “thói” nhân cơ hội ra ngoài nước để tuyên bố này nọ. Anh nói, anh có thể sẵn sàng trả lời phỏng vấn khi anh ở trong nước hơn.

(Cũng nhân lần ấy, anh đã trao cho họa sĩ, nhà văn Khánh Trường một bài thơ anh viết nhân ghé qua một sân bay ở Đức để đăng trên tạp chí Hợp Lưu. Rất tiếc tôi không tìm ra bài thơ “Chiều Hamburg” ấy khi viết những dòng này).

*
Một điều tôi không hề thấy ở anh: tham vọng cá nhân về quyền lợi hay danh vị. Người ta có thể nói, với những đóng góp và thành đạt của anh thì chuyện danh vị anh không cần quan tâm, nó vẫn đến. Tôi lại không nghĩ đơn giản như vậy. Có những thứ, với ai đó, bao nhiêu cũng chưa vừa, chưa đủ. Qua bao nhiêu lần có dịp trao đổi với anh tôi chưa một lần nhận thấy danh vị là mối quan tâm của anh. Đó là điều đáng nói. Về quyền lợi thì có lẽ cuộc sống thật của anh đã đủ để cho câu trả lời.

Hôm nay anh đã đi xa rồi. Không có cơ hội về tận nơi để chào tiễn biệt anh, tôi chỉ thầm cầu mong anh ra đi phần nào thanh thản. Tôi nghĩ đến các cháu, con anh chị. Những cô cậu bé ngày nào năm xưa tôi đã từng gặp giờ đã là những cá nhân thành đạt xứng đáng trong các chọn lựa riêng. Dẫu biết rằng, thành quả của con cái là do bản lĩnh và nỗ lực tự thân của con cái, nhưng dù sao thì bố mẹ cũng có quyền tự hào về những thành quả ấy. Và tôi vững tin rằng, trước khi ra đi anh cũng có thể vui lòng và tự hào về những thành đạt và đóng góp của các cháu. Tôi tự nhủ, đó phải là niềm vui, niềm hạnh phúc không nhỏ anh đã có được và mang theo.

Tôi đã viết “phần nào”, với chút xót xa chủ quan. Vì tôi vẫn luôn đau đáu nghĩ về những bậc cha chú, những bậc đàn anh đã suốt đời tận hiến tim óc cho mong ước về một quê hương tươi đẹp, hạnh phúc hơn. Ra đi mà được nhìn thấy thế giới tốt đẹp hơn khi mình đến thì, đã là người tâm huyết với quê nhà, ai chẳng ước mơ. Và riêng anh, anh đã tích cực đem tim óc mình để đóng góp vào đó. Anh đã từng mong những can gián, những đề nghị của anh được lắng nghe, được hiểu, được bàn bạc, đánh giá tương xứng và những gì phù hợp được biến thành những thành quả thiết thực, có lợi cho quê hương. Tiếc thay, khi ra đi anh vẫn chưa được nhìn thấy một quê hương đang đổi thay như anh hằng mong ước.


Cảm nghĩ ấy của tôi thực ra không phải hôm nay mới quay về. Buổi trưa ngồi nhìn anh ở quán ăn nọ tôi đã cảm thấy dâng lên nỗi xót xa ấy. Anh đi, anh còn để lại rất nhiều việc phải làm, để làm. Thôi thì phiên gác đã đổi, đó là luật trời. Xin anh yên nghỉ.

